% GENERATED FILE. Edit canonical lessons, then run npm run generate:books.
% canonical-source-hash: fnv1a64:d3af54ba676aae74
% canonical-lessons: FA-C463-sowghati, FA-C463-sanaye-e-dasti, FA-C463-gelim, FA-C463-javaher, FA-C463-atr

\chapter{Souvenirs}
\label{ch:fa-souvenirs}
\hlchaptermodality{463}

\begin{chapteropening}
\textbf{By the end of this chapter:} \emph{I can buy presents to take home: a souvenir, handicrafts, a kilim, jewellery and perfume.}

Complete the last lesson of chapter 463: I can buy presents to take home: a souvenir, handicrafts, a kilim, jewellery and perfume.
\end{chapteropening}

\section[\texorpdfstring{\fa{سوغاتی} (sowghâti)}{سوغاتی (sowghâti)}]{\fa{سوغاتی} (sowghâti) — a souvenir, a present brought back from a trip}

\label{lesson:FA-C463-sowghati}



\begin{quote}
Before the new one: say the Persian for to sell, then the Persian for a box.
\end{quote}

\subsection*{You'll want to know: \fa{سوغاتی}}
\textbf{\fa{سوغاتی}} (\emph{sowghâti}) — \textquotedblleft{}a souvenir, a present brought back from a trip\textquotedblright{}.

\begin{quote}
\fa{از} \fa{سفر} \fa{برای} \fa{مادرم} \fa{سوغاتی} \fa{آوردم}.
\end{quote}

\begin{quote}
\emph{Az safar barâ-ye mâdaram sowghâti âvardam.}
\end{quote}

\begin{quote}
I brought my mother a present back from the trip.
\end{quote}

Bringing \emph{sowghâti} for family and friends is expected after a trip.

\begin{cousinweb}
\emph{sowghâti} is built from \emph{sowghât} 'a gift from a journey', usually traced to a Turkic word for a present, and the Persian ending \emph{-i}. EtymDB does not list it.
\end{cousinweb}

\begin{grammarlens}[title={What you've built}]
One more word for telling, explaining and describing at B1.
\end{grammarlens}

\subsection*{Your turn}
\begin{itemize}
\raggedright
  \item \emph{Say it:} \emph{sowghâti}
  \item \emph{Say it:} \emph{sowghâti}, once more
  \item \emph{Say it:} a sentence of your own with \emph{sowghâti}
  \item \emph{Recall:} say the Persian for to sell, then the Persian for a box, then say \emph{sowghâti} again
\end{itemize}

\begin{tcolorbox}[breakable,colback=teal!4,colframe=teal!35!black,title={Before you move on}]
What does \fa{سوغاتی} mean? (\textquotedblleft{}A souvenir, a present brought back from a trip\textquotedblright{}.) Say it once more.
\end{tcolorbox}
\section[\texorpdfstring{\fa{صنایع} \fa{دستی} (sanâye'-e dasti)}{صنایع دستی (sanâye'-e dasti)}]{\fa{صنایع} \fa{دستی} (sanâye'-e dasti) — handicrafts, crafts}

\label{lesson:FA-C463-sanaye-e-dasti}



\begin{quote}
Before the new one: say the Persian for a box, then the Persian for a souvenir.
\end{quote}

\subsection*{You'll want to know: \fa{صنایع} \fa{دستی}}
\textbf{\fa{صنایع} \fa{دستی}} (\emph{sanâye'-e dasti}) — \textquotedblleft{}handicrafts, crafts\textquotedblright{}.

\begin{quote}
\fa{در} \fa{بازار} \fa{صنایع} \fa{دستی} \fa{خریدیم}.
\end{quote}

\begin{quote}
\emph{Dar bâzâr sanâye'-e dasti kharidim.}
\end{quote}

\begin{quote}
We bought handicrafts in the bazaar.
\end{quote}

\emph{sanâye'} is an Arabic broken plural of \emph{san'at} 'craft, industry'; \emph{dasti} 'by hand'.

\begin{cousinweb}
\emph{sanâye'-e dasti} is built from \emph{sanâye'} 'crafts', the Arabic plural \emph{ṣanāʾiʿ} that EtymDB lists, and native \emph{dasti} 'by hand', from \emph{dast} 'hand'.
\end{cousinweb}

\begin{grammarlens}[title={What you've built}]
One more word for telling, explaining and describing at B1.
\end{grammarlens}

\subsection*{Your turn}
\begin{itemize}
\raggedright
  \item \emph{Say it:} \emph{sanâye'-e dasti}
  \item \emph{Say it:} \emph{sanâye'-e dasti}, once more
  \item \emph{Say it:} a sentence of your own with \emph{sanâye'-e dasti}
  \item \emph{Recall:} say the Persian for a box, then the Persian for a souvenir, then say \emph{sanâye'-e dasti} again
\end{itemize}

\begin{tcolorbox}[breakable,colback=teal!4,colframe=teal!35!black,title={Before you move on}]
What does \fa{صنایع} \fa{دستی} mean? (\textquotedblleft{}Handicrafts, crafts\textquotedblright{}.) Say it once more.
\end{tcolorbox}
\section[\texorpdfstring{\fa{گلیم} (gelim)}{گلیم (gelim)}]{\fa{گلیم} (gelim) — a kilim, a flat-woven rug}

\label{lesson:FA-C463-gelim}



\begin{quote}
Before the new one: say the Persian for a souvenir, then the Persian for handicrafts.
\end{quote}

\subsection*{You'll want to know: \fa{گلیم}}
\textbf{\fa{گلیم}} (\emph{gelim}) — \textquotedblleft{}a kilim, a flat-woven rug\textquotedblright{}.

\begin{quote}
\fa{روی} \fa{زمین} \fa{اتاق} \fa{یک} \fa{گلیم} \fa{قرمز} \fa{بود}.
\end{quote}

\begin{quote}
\emph{Ru-ye zamin-e otâq yek gelim-e qermez bud.}
\end{quote}

\begin{quote}
There was a red kilim on the floor of the room.
\end{quote}

A \emph{gelim} has no pile, unlike a \emph{farsh} or \emph{qâli}.

\begin{cousinweb}
\emph{gelim} is a Persian word that EtymDB links with Greek \emph{kálymma} 'a cover, a covering'.
\end{cousinweb}

\begin{grammarlens}[title={What you've built}]
One more word for telling, explaining and describing at B1.
\end{grammarlens}

\subsection*{Your turn}
\begin{itemize}
\raggedright
  \item \emph{Say it:} \emph{gelim}
  \item \emph{Say it:} \emph{gelim}, once more
  \item \emph{Say it:} a sentence of your own with \emph{gelim}
  \item \emph{Recall:} say the Persian for a souvenir, then the Persian for handicrafts, then say \emph{gelim} again
\end{itemize}

\begin{tcolorbox}[breakable,colback=teal!4,colframe=teal!35!black,title={Before you move on}]
What does \fa{گلیم} mean? (\textquotedblleft{}A kilim, a flat-woven rug\textquotedblright{}.) Say it once more.
\end{tcolorbox}
\section[\texorpdfstring{\fa{جواهر} (javâher)}{جواهر (javâher)}]{\fa{جواهر} (javâher) — jewels, jewellery}

\label{lesson:FA-C463-javaher}



\begin{quote}
Before the new one: say the Persian for handicrafts, then the Persian for a kilim.
\end{quote}

\subsection*{You'll want to know: \fa{جواهر}}
\textbf{\fa{جواهر}} (\emph{javâher}) — \textquotedblleft{}jewels, jewellery\textquotedblright{}.

\begin{quote}
\fa{در} \fa{موزه} \fa{جواهر} \fa{قدیمی} \fa{دیدیم}.
\end{quote}

\begin{quote}
\emph{Dar muze javâher-e qadimi didim.}
\end{quote}

\begin{quote}
We saw old jewels in the museum.
\end{quote}

An Arabic broken plural of \emph{jowhar}; a jeweller is a \emph{javâherforush}.

\begin{cousinweb}
\emph{javâher} came back into Persian from Arabic \emph{jawāhir} 'jewels' (EtymDB), which itself was taken from Persian \emph{gowhar} 'jewel', from Middle Persian and Old Persian \emph{gauthra} 'substance, essence'.
\end{cousinweb}

\begin{grammarlens}[title={What you've built}]
One more word for telling, explaining and describing at B1.
\end{grammarlens}

\subsection*{Your turn}
\begin{itemize}
\raggedright
  \item \emph{Say it:} \emph{javâher}
  \item \emph{Say it:} \emph{javâher}, once more
  \item \emph{Say it:} a sentence of your own with \emph{javâher}
  \item \emph{Recall:} say the Persian for handicrafts, then the Persian for a kilim, then say \emph{javâher} again
\end{itemize}

\begin{tcolorbox}[breakable,colback=teal!4,colframe=teal!35!black,title={Before you move on}]
What does \fa{جواهر} mean? (\textquotedblleft{}Jewels, jewellery\textquotedblright{}.) Say it once more.
\end{tcolorbox}
\section[\texorpdfstring{\fa{عطر} ('atr)}{عطر ('atr)}]{\fa{عطر} ('atr) — a perfume, a scent}

\label{lesson:FA-C463-atr}



\begin{quote}
Before the new one: say the Persian for a kilim, then the Persian for jewels.
\end{quote}

\subsection*{You'll want to know: \fa{عطر}}
\textbf{\fa{عطر}} (\emph{'atr}) — \textquotedblleft{}a perfume, a scent\textquotedblright{}.

\begin{quote}
\fa{خواهرم} \fa{برای} \fa{تولدم} \fa{یک} \fa{عطر} \fa{خرید}.
\end{quote}

\begin{quote}
\emph{Khâharam barâ-ye tavallodam yek 'atr kharid.}
\end{quote}

\begin{quote}
My sister bought me a perfume for my birthday.
\end{quote}

\emph{'Atr-e gol-e mohammadi} — rose scent, made from the damask rose.

\begin{cousinweb}
\emph{'atr} was borrowed from Arabic \emph{ʿiṭr} 'perfume' (EtymDB).
\end{cousinweb}

\begin{grammarlens}[title={What you've built}]
One more word for telling, explaining and describing at B1.
\end{grammarlens}

\subsection*{Your turn}
\begin{itemize}
\raggedright
  \item \emph{Say it:} \emph{'atr}
  \item \emph{Say it:} \emph{'atr}, once more
  \item \emph{Say it:} a sentence of your own with \emph{'atr}
  \item \emph{Recall:} say the Persian for a kilim, then the Persian for jewels, then say \emph{'atr} again
\end{itemize}

\begin{tcolorbox}[breakable,colback=teal!4,colframe=teal!35!black,title={Before you move on}]
What does \fa{عطر} mean? (\textquotedblleft{}A perfume, a scent\textquotedblright{}.) Say it once more.
\end{tcolorbox}
